\chapter{Discussion} \label{Chap5}

\section{A Possible Mechanistic Explanation of GRPO's Robustness}

The core mechanism of GRPO is group-relative advantage: for $K$ rollouts from the same initial state, the advantage function computes an individual's relative performance against the within-group mean rather than the absolute reward value. The observations of Chapter~4 show that GRPO remains stable under both reward conditions, whereas PPO reward-hacks under the Misaligned condition and collapses or stalls under the Sparse condition. We explain this difference at the mechanistic level below.

Under the Misaligned Reward condition, PPO's GAE mechanism may form a positive-feedback loop: the policy accidentally discovers that pushing and pulling the door increases the door-openness reward; GAE encodes this absolute increase as a positive gradient signal; the updated policy becomes better at pushing the door, the reward grows further, and the gradient becomes stronger, driving the behaviour to be even more extreme. GRPO's group-relative mechanism may instead break this loop: when most samples in a group obtain a high door-openness reward, the within-group mean rises together and the individual's relative advantage approaches zero; even a high absolute reward produces no strong gradient. Only behaviour that genuinely outperforms the other samples in the group is reinforced. Because the passing-the-door reward is extremely hard to obtain, no sample can stand out on that dimension, so the policy does not over-favour any single component; this is consistent with the observation in Section~4.2.1 that GRPO does not amplify the door-openness component.

Under the Sparse Reward condition, PPO faces a double dilemma: the reward-hacking behaviour learned under the Misaligned condition no longer yields any return, and the sparse binary signal provides too little information to guide the policy out of its erroneous behavioural pattern. Under sparse rewards its GAE advantage estimate is either too noisy (causing deterministic collapse) or too weak (causing gradient vanishing), and the common root of the two failure modes observed in Section~4.2.2 is the lack of a stable signal scale. In contrast, GRPO's within-group normalization always scales the advantage signal to a stable magnitude independent of the absolute scale of the external reward; since the policy is not locked into hacking behaviour under the Misaligned condition, it does not need to recover and can adapt to the Sparse condition directly, as shown by the cross-condition comparison in Section~4.1.

In addition, GRPO's log-$\sigma$ remains conservative under the Sparse condition while PPO's exploration level diverges strongly across seeds (Section~4.2.2). This suggests that the conservativeness of log-$\sigma$ may be a key factor enabling GRPO to avoid both deterministic collapse and gradient vanishing under extremely sparse rewards, although a causal verification of this hypothesis is left to future work.

\section{Limitations}

\textbf{Weak baseline.} The baseline's 5\% success rate is low and may limit the upper bound of fine-tuning performance. Future work could use a stronger baseline (e.g., a model pretrained with imitation learning or more rounds of curriculum training) to verify the generality of our conclusions under a stronger baseline.

\textbf{State-distribution bias in KL divergence.} Our KL-divergence analysis is based on 1000 observations collected under the baseline policy's rollout distribution. If a fine-tuned policy drives the robot into states the baseline never visited (e.g., PPO's extreme door-manipulation behaviour), the KL divergence may be systematically underestimated, because the policy differences on those out-of-distribution states are not sampled. In practice this limitation means the true KL difference between PPO and GRPO may be larger than reported here. Future work should use a mixture of rollouts from multiple policies to cover a more complete state distribution.

\textbf{Hyperparameter sensitivity of GRPO.} Limited by single-GPU memory, we use only one set of GRPO hyperparameters: group size $K=4$ and learning rate $3\times10^{-5}$. The sensitivity of GRPO to the group size remains to be studied; a larger $K$ may enhance the statistical stability of within-group normalization, but may also weaken the discriminative power of the relative advantage.

\textbf{Limited verification of the mechanistic explanation.} The KL ablation in Section~4.3 rules out the influence of the ``absence of the KL penalty'' design choice on stability, but PPO and GRPO also differ in the presence of a critic, and the causal role of group-relative advantage has not been fully separated by a dedicated control experiment. Future work could further verify this via a critic-free PPO or a critic-based GRPO control.
